
\subsection{3.1, Matrix transformations}

\begin{problems}
%a
%
\problem % Dan's notes
True or false: If $A$ is a matrix and $T$ is the associated matrix 
transformation, then the statement that $Ax = b$ is consistent is 
equivalent to the statement that $b$ is in the range of $T$.

\begin{solution}
  True. This follows directly from the definition of range and the fact that $T$
  is the transformation $T(x) = Ax$: \newline %
  The matrix equation $Ax = b$ is consistent $\iff$ $T(x) = b$ is consistent
    $\iff$ $b$ is in the range of $T$.
\end{solution}

\problem % Dan's notes
There is a matrix $A$ so that the domain of the associated matrix transformation
$T$ is a line in ${\bf R}^3$.

\begin{solution}
  False: the domain of $T$ is (all of) ${\bf R}^n$, where $n$ is the number of
  columns of $A$.  If $A$ has one column, then the domain of $T$ is ${\bf
    R}$. If $A$ has three columns, then the domain of $T$ is ${\bf R}^3$.  The
  domain of $T$ can never be a line in ${\bf R}^3$.
\end{solution}


\problem % Dan's notes
What does the matrix $\mat{-1 0; 0 -1}$ do to vectors in ${\bf R}^2$?

\begin{solution}
  $\mat{-1 0; 0 -1}\vec{x y} = \vec{-x -y} = - \vec{x y}$, so the matrix sends
  each vector $v$ to $-v$.  In other words, the matrix implements scalar
  multiplication by $-1$.
\end{solution}



\problem % Dan's notes
What does the matrix $\mat{1 0 0; 0 0 0; 0 0 0}$ do to vectors in ${\bf R}^3$?

\begin{solution}
It naturally projects each vector onto the $x$-axis by turning its second and third coordinates into zeros:
\[\mat{1 0 0; 0 0 0; 0 0 0} \vec{x y z} = \vec{x 0 0}.\]
\end{solution}

\problem % Dan's notes
What does the matrix $\mat{0 0; 1 0; 0 1}$ do to vectors in ${\bf R}^2$?

\begin{solution}
The matrix takes vectors from ${\bf R}^2$ and naturally inserts them into the $yz$-plane in ${\bf R}^3$.  If $\vec{x_1 x_2}$ is any vector in ${\bf R}^2$, then
\[\mat{0 0; 1 0; 0 1} \vec{x_1 x_2} = \vec{0 x_1 x_2}.\]
\end{solution}


\problem % Created for ILA
Let $T$ be the transformation $T\vec{x y} = \mat{0 -1; 1 0} \vec{x y}$.
Describe $T$ geometrically.

\begin{solution}
$T$ rotates vectors in ${\bf R}^2$ by $\frac{\pi}{2}$ radians counterclockwise. 
\end{solution}

\problem % Created for ILA
Let $T$ be the transformation $T\vec{x y} = \mat{0 1; 1 0} \vec{x y}$.  Describe
$T$ geometrically.

\begin{solution}
Rotation across the line $y=x$.  We can check by
  computing $T \vec{x y} = \vec{y x}$.
\end{solution}


\problem % Created for ILA
Fill in the blanks: If $T: {\bf R}^{5} \to {\bf R}^3$ is a matrix
transformation $T(x) = Ax$, then $A$ has $\underline{\qquad}$
rows and $\underline{\qquad}$ columns.

\begin{solution}
  $A$ has $3$ rows and $5$ columns.  This may seem counterintuitive,
  but we can check: for $T(x)$ to make sense, $x$ must be in ${\bf R}^5$,
  so $A$ has $5$ columns.  Moreover, $T(x)$ is in ${\bf R}^3$ so $A$
  must have $3$ rows.
\end{solution}


\problem % Created for ILA
Let $A$ be a $7 \times 6$ matrix, and define a matrix transformation
by $T(x) = Ax$.  What is the domain of $T$?  What is the codomain of $T$?

\begin{solution}
  The domain of $T$ is ${\bf R}^6$ and the codomain is ${\bf R}^7$.
  We can check: for $x$ to be in the domain of $T$, we need $Ax$ to
  make sense, so $x$ must be in ${\bf R}^6$.  The vector $Ax$ lives in
  the column space of $A$ (which is the codomain of $T$), namely
  ${\bf R}^7$.

  In general, For a $k \times \ell$ matrix $A$, the transformation
  $T(x) = Ax$ has domain ${\bf R}^{\ell}$ and codomain ${\bf R}^k$.
\end{solution}

\problem Write a $2 \times 2$ matrix $A$ so that the transformation
$T\vec{x y} = A\vec{x y}$ reflects vectors across the $x$-axis.

% Created for ILA
\begin{solution}
$A = \mat{1 0; 0 -1}$.  We can check by verifying that 
$A \vec{x y} = \vec{x -y}$.    
\end{solution}


\problem Write a $2 \times 2$ matrix $A$ so that the transformation
$T\vec{x y} = A\vec{x y}$ stretches all vectors by a factor of $4$.

% Created for ILA
\begin{solution}
$A = \mat{4 0; 0 4}$.  We can check by verifying 
$A \vec{x y} = \vec{4x 4y} = 4 \vec{x y}$.
\end{solution}

\problem Write a $2 \times 2$ matrix $A$ so that the transformation
$T\vec{x y} = A\vec{x y}$ reflects vectors across the $y$-axis.

% Created for ILA
\begin{solution}
$A = \mat{-1 0; 0 1}$.  We can check by verifying that
$A \vec{x y} = \vec{-x y}$.
\end{solution}

\problem Write a $2 \times 2$ matrix $A$ so that the transformation
$T\vec{x y} = A\vec{x y}$ stretches vectors along the $x$-axis by a factor of
$3$, while leaving the $y$-axis fixed.%

% Created for ILA
\begin{solution}
$A = \mat{4 0; 0 4}$.  We can check by verifying 
$A \vec{x y} = \vec{4x 4y} = 4 \vec{x y}$.
\end{solution}


\problem Consider the transformation
$T(x) = \mat{1 -1 9; 0 2 4; 0 1 2; 0 1 2} x$.
\begin{subproblems}
  \subproblem What is the domain of $T$?
  \subproblem What is the codomain of $T$?
  \subproblem What is the range of $T$?
\end{subproblems}

\begin{solution} $T$ is the matrix transformation $T(x) = Ax$ for the
  $4 \times 3$ matrix $A = \mat{1 -1 9; 0 2 4; 0 1 2; 0 1 2}$.
  \begin{subproblems}
    \subproblem The domain of $T$ is ${\bf R^3}$.  %
    \subproblem The codomain of $T$ is ${\bf R}^4$. %
    \subproblem The range of $T$ is the column space of $A$.  As we
    know from the the past, the pivot columns
    \[A = \mat{1 -1 9; 0 2 4; 0 1 2; 0 1 2} \xrightarrow[]{RREF}
      \mat{1 0 11; 0 1 2; 0 0 0; 0 0 0}.\] Thus, the first two columns
    are pivot columns, so range$(T)$ is
    \[\textup{Span}\left\{\vec{1 0 0 0}, \vec{-1 2 1 1}\right\}.\]
  \end{subproblems}
\end{solution}


\problem Consider the transformation
\[T(x) = \mat{2 2 3 1 2; 3 4 0 0 -1}x.\]
\begin{subproblems}
  \subproblem What is the domain of $T$?
  \subproblem What is the codomain of $T$?
  \subproblem What is the range of $T$?
\end{subproblems}

\begin{solution} $T(x) = Ax$ for the $2 \times 5$ matrix
  $A = \mat{2 2 3 1 2; 3 4 0 0 -1}$.
  \begin{subproblems}
    \subproblem The domain of $T$ is ${\bf R^5}$.  %
    \subproblem The codomain of $T$ is ${\bf R}^2$. %
    \subproblem The range of $T$ is the column space of $A$. We could
    do one step of row-reduction to see that the first two columns of
    $A$ are pivot columns.  Therefore, the first two columns are
    linearly independent vectors in ${\bf R}^2$, so the range of $T$
    is ${\bf R}^2$.
  \end{subproblems}
\end{solution}



\problem Is there a matrix transformation $T(x) = Ax$ that satisfies
the following?
\[T\vec{1 3} = \vec{1 0}, \qquad T \vec{2 6} = \vec{-1 0}.\] (hint:
recall the matrix multiplication properties $A(u+v) = Au+Av$,
$A(cu) = cAu$ for vectors $u$ and $v$ and scalars $c$).

\begin{solution}
  No.  If $T$ is a matrix transformation $T(x) = Ax$ for some matrix
  $A$, then from the property $A(2u) = 2 A(u)$ it must be true that
  \[A \vec{2 6} = 2 A \vec{1 3},\] or in other words
  $T \vec{2 6} = 2 T\vec{1 3}$.

  However, $T\vec{2 6} = \vec{-1 0}$ but $2 T \vec{1 3} = \vec{2 0}$,
  so $T \vec{2 6} \neq 2 T\vec{1 3}$.
\end{solution}

\end{problems}\setcounter{problem}{0}

\newpage

\begin{center} \underline{3.2, One-to-one and onto
    transformations} \end{center}

\begin{problems}
  \problem % Dan's notes
  It is possible for the matrix transformation associated to a $5 \times 6$
  matrix to be both one-to-one and onto.

\begin{solution}
  False.  A matrix cannot have more pivots than rows or columns, so the matrix
  has a max of $5$ pivots.  This means it is possible for the matrix to have a
  pivot in every row, so the corresponding transformation might be onto.
  However, the matrix will have at least one pivot without a column, so the
  corresponding transformation $T$ cannot be one-to-one.
\end{solution}


\problem % Dan's notes
The transformation $T: {\bf R}^3 \to {\bf R}^3$ given by projection onto the
$yz$-plane is onto.

\begin{solution}
False.  Recall that in order for the transformation to be onto, its range must be all of ${\bf R}^3$.  However, the range of $T$ is the $yz$-plane.  In other words, the range is $\Span\left\{\vec{0 1 0}, \vec{0 0 1}\right\}$, which is only a $2$-dimensional subspace of ${\bf R}^3$, not all of ${\bf R}^3$.
\end{solution}


\problem % Dan's notes
The transformation $T: {\bf R}^2 \to {\bf R}^2$ given by rotation by $\pi/5$
radians clockwise is onto.

\begin{solution}
  True.  Let $y$ be \textit{any} vector in ${\bf R}^2$.  Then for the vector $x$
  obtained by rotating $y$ by $\pi/5$ radians \textit{counterclockwise}, we get
  $T(x) = y$.  This shows that every vector in ${\bf R}^2$ is in the range of
  $T$.  In other words, range$(T) = {\bf R}^2$, so $T$ is onto.
\end{solution}

\problem %
If a matrix $A$ has more rows than columns, then the matrix transformation $T$
given by $T(x)=Ax$ is not onto.

\begin{solution}
True.  Since $A$ cannot have more pivots than columns,
    and since there are more rows than columns, we conclude $A$ cannot
    have a pivot in every row.  
\end{solution}

\problem % 
Suppose $A = \mat{1 2 3; 4 5 6; 3 6 9}$. Then the transformation $T(x) = Ax$ is one-to-one.

\begin{solution}
False:
    $\mat{1 2 3; 4 5 6; 3 6 9} \xrightarrow[R_3 = R_3 - 3R_1]{R_2=
      R_2-4R_1} \mat{1 2 3; 0 -3 -6; 0 0 0}$, so $A$ does not have a
    pivot in every column and therefore $T$ is not one-to-one.
\end{solution}

\problem True or false: Suppose $A$ is
a matrix and $A \vec{-3 5 8 1}$ is the zero vector.  Then the transformation
$T(x) = Ax$ cannot be one-to-one.

\begin{solution}
  True.  Since $A \vec{-3 5 8 1}$ is the zero vector, the homogeneous
  system $Ax = 0$ has more than just the trivial solution.  Therefore, 
  there are infinitely many vectors $x$ so that $Ax = 0$, so $T(x)=0$
  has infinitely many solutions and $T$ is not one-to-one.
\end{solution}

\problem % Dan's notes
Is there an onto matrix transformation $T(x) = Ax$ with domain ${\bf R}^2$ and
codomain ${\bf R}^3$?  If so, write one down.  If not, explain why not.

\begin{solution}
  No.  In order for $T$ to have domain ${\bf R}^2$ and codomain ${\bf R}^3$, the
  matrix $A$ must be $3 \times 2$, so $A$ has a max of $2$ pivot columns.
  Therefore, the column space of $A$ (which is the range of $T$) is at most a
  $2$-dimensional subspace of ${\bf R}^3$, so $T$ cannot be onto.
\end{solution}

\problem % Dan's notes
Is there a one-to-one matrix transformation with domain ${\bf R}^2$ and codomain
${\bf R}^3$?  If so, write one down.  If not, explain why not.

\begin{solution}
Yes.  Some examples of $3 \times 2$ matrice with a pivot in each column are
\[A = \mat{1 0; 0 1; 0 0}, \qquad A = \mat{1 1; 0 1; 0 0}, 
\qquad A = \mat{0 0; 5 1; 1 0}.\]
\end{solution}



\problem Which of the following transformations are one-to-one?
\begin{subproblems}
  \subproblem $T: {\bf R}^2 \to {\bf R}^2$ that naturally projects
  vectors onto the $x$-axis. %
  \subproblem $T(x, y, z) = (x, y+z)$.
\end{subproblems}

\begin{solution}
  \begin{subproblems}
    \subproblem No: for example, $T(0, 1) = (0, 0)$ and
    $T(0, 0) = (0, 0)$. Alternatively, we could write the matrix for
    $T$ (which we learned in section 3.1) and observe that it has $2$
    columns but only one pivot.%
    \subproblem No: for example, $T(1, 0, 0) = (1, 0)$ and
    $T(1, 1, -1) = (1, 0)$. %
  \end{subproblems}
\end{solution}

\problem Either write a $2 \times 3$ matrix $A$ that so that the
matrix transformation $T(x)=Ax$ is one-to-one, or justify why no such
matrix $A$ exists.

\begin{solution}
  Any $2 \times 3$ matrix $A$ has $3$ columns but at most $2$ pivots,
  so it cannot have a pivot in every column, hence $T(x) = Ax$ cannot
  be one-to-one.
\end{solution}

\problem % Made for ILA
We are told that a matrix $A$ has four columns and three pivots, and let $T$ be the associated linear transformation.
\begin{subproblems}
  \subproblem Geometrically, what can we conclude about the range of
  $T$?  Do we know in which ${\bf R}^k$ we can find the range of $T$?
  Is it possible that $T$ is onto? %
  \subproblem Geometrically, what can we conclude about the set of all
  solutions to the equation $T(x) = 0$?  Do we know in which
  ${\bf R}^\ell$ we can find this solution set?  Is it possible that
  $T$ is one-to-one?
\end{subproblems}

\begin{solution}
  \begin{subproblems}
    \subproblem The range of $T$ is the column span of $A$.  Since $A$
    has three pivots, this means that the range of $T$ has dimension
    $3$.  We do not know where the range of $T$ lives, because we were
    not told the number of rows of $A$ and we cannot tell from the
    information given.
    %
    \subproblem The solution set for $T(x) = 0$ is just the null space
    of $A$.  Since $A$ has $3$ pivots and $4$ columns, the system
    $Ax=0$ will have one free variable and thus infinitely many
    solutions.  Therefore, $T$ is not one-to-one and the solution set
    to $T(x)=0$ is a one-dimensional subspace of ${\bf R}^4$.
  \end{subproblems}
\end{solution}


\problem Either write a $3 \times 4$ matrix $A$ so that the matrix
transformation $T(x)=Ax$ is onto, or justify why no such matrix $A$
exists.

\begin{solution}
  Many examples possible, for example
  $A = \mat{1 0 0 0; 0 1 0 0; 0 0 1 0}$.  We can see directly that
  every vector $\vec{x_1 x_2 x_3}$ in ${\bf R}^3$ is in the range of
  the transformation $T(x)=Ax$, since
  \[A\vec{x_1 x_2 x_3 x_4} = \mat{1 0 0 0; 0 1 0 0; 0 0 1 0}
    \vec{x_1 x_2 x_3 0} = \vec{x_1 x_2 x_3}.\]
\end{solution}


\problem Write a linear transformation $T$ with domain ${\bf R}^3$ and
codomain ${\bf R}^2$ whose range is the line $x_2 = -2x_1$.

\begin{solution}
  We need to find a $2 \times 3$ matrix $A$ whose column space is the
  line $x_2 = -2x_1$.  In other words
  $\Col(A) = \textup{Span}\left\{ \vec{1 -2}\right\}$.  We need $A$ to
  have exactly one pivot column, and each column must be a scalar
  multiple of $\vec{1 -2}$.  Many examples are possible, for example:
  \[A = \mat{1 0 0; -2 0 0}, \qquad A = \mat{1 2 3; -2 -4 -6}.\]
\end{solution}

\end{problems}\setcounter{problem}{0}
\newpage

\subsection{3.3, Linear transformations}

\begin{problems}
%a
%
\problem % Dan's notes
Define $T: {\bf R} \to {\bf R}$ by $T(x) = x+1$.  Is $T$ a linear transformation?

\begin{solution}
No.  Recall that every linear transformation must send the origin to the origin.  However, $T(0) = 1$, so $T$ cannot be linear.  We could additionally check that $T$ fails other linearity properties.  For example, $T(1 + 2) \neq T(1) + T(2)$.
\end{solution}

\problem % Dan's notes
Suppose $T: {\bf R}^2 \to {\bf R}^3$ is a linear transformation satisfying
\[T\vec{1 1} = \vec{3 3 1} \quad \textup{ and } \quad T \vec{2 1} = \vec{3 1 1}.\]
What is $T \vec{1 0}$?

\begin{solution}
First we write $\vec{1 0}$ as a linear combination of $u = \vec{1 1}$ and $v = \vec{2 1}$.
\[x_1 u + x_2 v = \vec{1 0}, \qquad x_1 \vec{1 1} + x_2 \vec{2 1} = \vec{1 0}.\]
This gives $x_1 = -1$ and $x_2 = 1$, so 
\[-u + v= \vec{1 0}.\]
By linearity of $T$:
\[T\vec{1 0} = T(-u+v) = T(-u)+T(v) = -T(u)+T(v)
= -\vec{3 3 1} + \vec{3 1 1} = \vec{0 -2 0}.\]
% The matrix used for $T$ is $A = \mat{0 3; -2 5; 0 1}$
\end{solution}

\problem 
Suppose $T: {\bf R}^2 \to {\bf R}^2$ is a linear transformation satisfying
\[T\vec{3 -1} = \vec{1 1} \quad \textup{ and } \quad T \vec{1 2} = \vec{4 4}.\]
Find the standard matrix $A$ for $T$, and use it to find $T\vec{11 1}$.

\begin{solution}
  We know that $A = \mat{T\vec{1 0} T\vec{0 1}}$, and we use linearity of $T$ to
  find both columns of $A$.  Let $u = \vec{3 -1}$ and $v = \vec{1 2}$.  Work
  using the basics of chapter 1 gives us
\[\vec{1 0} = \dfrac{2}{7} u + \dfrac{1}{7} v \quad
\textup{ and } \quad 
\vec{0 1} = \dfrac{-1}{7} u + \dfrac{3}{7} v.
\]
Therefore, 
\[T \vec{1 0} = T\left(\dfrac{2}{7}u + \dfrac{1}{7} v\right)
= \dfrac{2}{7} T(u) + \dfrac{1}{7} T(v) 
= \dfrac{2}{7} \vec{1 1} + \dfrac{1}{7} \vec{4 4}
= \vec{6/7 6/7}\]
and
\[T \vec{0 1} = T\left(\dfrac{-1}{7}u + \dfrac{3}{7} v\right)
= \dfrac{-1}{7} T(u) + \dfrac{3}{7} T(v) 
= \dfrac{-1}{7} \vec{1 1} + \dfrac{3}{7} \vec{4 4}
= \vec{11/7 11/7.}\]
Therefore, 
\[A = \mat{T\vec{1 0} T\vec{0 1}} = \mat{6/7 11/7; 6/7 11/7}.\]
To finsih, we compute:
\[T\vec{11 1} = A \vec{11 1} = 
\mat{6/7 11/7; 6/7 11/7} \vec{11 1} = \vec{11 11}.\]
If we had not been asked to compute $A$, we could have found $T\vec{11 1}$ directly by using linearity of $T$, by writing $\vec{11 1}$ as the linear combination of $u$ and $v$.  We would get
$\vec{11 1} = 3u + 2v$, so 
\[T\vec{11 11} = T(3u+2v) = 3T(u)+2T(v) = 3 \vec{1 1} + 2 \vec{4 4}
= \vec{11 11}.\]
\end{solution}



\problem % Dan's notes
Find the matrix for the linear transformation $T: {\bf R}^3 \to {\bf R}^3$ that
rotates about the $z$-axis by $\pi$ radians and then scales by $2$.

\begin{solution}
  This is the composition $T \circ U$ where $T$ is the scaling operation and $U$
  is the rotation.  The matrix for $T$ is $\mat{2 0 0; 0 2 0; 0 0 2}$.  Note $U$
  keeps the $z$-axis fixed while rotating the $xy$ portion by $\pi$ radians, so
  the matrix for $U$ is $\mat{-1 0 0; 0 -1 0; 0 0 1}$.
  \[(T \circ U)(x) = \mat{2 0 0; 0 2 0; 0 0 2} \mat{-1 0 0; 0 -1 0; 0 0 1} =
    \mat{-2 0 0; 0 -2 0; 0 0 2}.\] It turns out that these two transformations
  happen to satisfy $T \circ U = U \circ T$, so we would have gotten the correct
  answer even if we had mixed up the order of operations.
\end{solution}

\problem % Dan's notes
Suppose $T: {\bf R}^3 \to {\bf R}^3$ is the function given by
\[T \vec{x y z} = \vec{z 0 x}.\] Is $T$ linear?  If so, what is the standard
matrix for $T$?

\begin{solution}
Yes, $T$ is linear, and we can verify that $T(x) = Ax$ for the matrix
\[A = \mat{T(e_1) T(e_2) T(e_3)} = \mat{0 0 1; 0 0 0; 1 0 0}.\]
We can check by computing
\[A \vec{x y z} = \mat{0 0 1; 0 0 0; 1 0 0}\vec{x y z}
= \vec{z 0 x}.\]
\end{solution}


\problem % Dan's notes
What does the matrix 
$\mat{-1/\sqrt{2} 1/\sqrt{2}; -1/\sqrt{2} -1/\sqrt{2}}$ 
do to vectors in ${\bf R}^2$?

\begin{solution}
  This matrix is $\mat{\cos(\theta) -\sin(\theta); \sin(\theta) \cos(\theta)}$
  for $\theta = 5\pi/4$ (or equivalently, $\theta = 225^\circ$).  Therefore,
  this matrix rotates each vector in ${\bf R}^2$ by $5\pi/4$ radians
  counterclockwise.
\end{solution}

\problem Is the identity transformation from ${\bf R}^n$ to ${\bf R}^n$ one-to-one?

\begin{solution}
  Yes: its standard matrix is the $n \times n$ identity matrix, which has a
  pivot in every row and column, so the identity transformation is both
  one-to-one and onto.
\end{solution}

\problem Which of the following transformations $T: {\bf R}^3
\to {\bf R}^3$ are linear?
\begin{subproblems}
  \subproblem $T(x_1, x_2, x_3) = (x_1-x_2, x_3 + 1, x_2)$ %
  \subproblem $T(x_1, x_2, x_3) = (x_1, |x_2|-x_1, x_3)$ %
\end{subproblems}

\begin{solution}
  \begin{subproblems}
    \subproblem No: $T(0, 0, 0) = (0, 1, 0)$.%
    \subproblem No: the absolute value term ruins the possibility
    that $T$ is linear.  If $u = (0, 1, 0)$ then
    \[T(u) = (0, 1, 0), \quad T(-u) = T(0,-1,0)= (0,1,0),\] so
    $T(-u) \neq -T(u)$. %
  \end{subproblems}
\end{solution}


\problem Which of the following transformations $T: {\bf R}^3
\to {\bf R}$ are linear?
\begin{subproblems}
\subproblem $T(a, b, c) = c-2b-a$
\subproblem $T(x, y, z) = xy-z$ %
\end{subproblems}

\begin{solution}
  \begin{subproblems}
    \subproblem Yes: $T(x) = Ax$ for $A = \mat{-1 -2 1}$. %
    \subproblem No: the product $xy$ ruins it.  For $u = \vec{1 1 0}$,
    we get
    \[T(u) = 1(1)-0 = 1, \qquad T(2u) = T\vec{2 2 0} = 2(2)-0 = 4.\]
    Therefore, $T(2u) \neq 2 T(u)$, so $T$ is not linear.
  \end{subproblems}
\end{solution}

\problem Suppose $T: {\bf R}^k \to {\bf R}^{\ell}$ is a one-to-one
linear transformation.  Which of the following are possible? Answer
yes or no in each case.
\begin{subproblems}
  \subproblem $k < \ell$
  \subproblem $k > \ell$
  \subproblem $k = \ell$.
\end{subproblems}

\begin{solution}
  Let $A$ be the standard matrix for $T$, so $A$ is an $\ell \times k$
  matrix.
  \begin{subproblems}
    \subproblem Yes.  If $k < \ell$ then $A$ has more rows than
    columns, so $A$ can have a pivot in every column, for example
    $A = \mat{1 0; 0 1; 0 0}$. %
    \subproblem No.  If $k > \ell$ then $A$ has more columns that
    rows, so $A$ cannot have a pivot in every column. %
    \subproblem Yes.  If $k = \ell$ then $A$ can have a pivot in every
    column, for example $A = \mat{1 0; 0 1}$.
  \end{subproblems}
\end{solution}

\problem Suppose $T: {\bf R}^a \to {\bf R}^b$ is a one-to-one
linear transformation.  Which of the following are possible? Answer
yes or no in each case.
\begin{subproblems}
  \subproblem $a < b$
  \subproblem $a > b$
  \subproblem $a = b$.
\end{subproblems}

\begin{solution}
  Let $A$ be the standard matrix for $T$, so $A$ is a $b \times a$
  matrix.
  \begin{subproblems}
    \subproblem No.  If $a < b$ then $A$ has more rows than columns,
    so $A$ cannot have a pivot in every row. %
    \subproblem Yes.  If $a > b$ then $A$ has more columns than rows,
    so $A$ can have a pivot in every row, for example if $a = 3$ and
    $b=2$ then $A$ is a $2 \times 3$ matrix.
    $A = \mat{1 0 0; 0 1 0}$. %
    \subproblem Yes.  If $a = b$ then $A$ can have a pivot in every
    row, for example $A = \mat{1 0; 0 1}$.
  \end{subproblems}
\end{solution}


\problem Consider the linear transformation
$T: {\bf R}^2 \to {\bf R}^4$ given by $T(a, b) = (a, b, b, a)$.
\begin{subproblems}
  \subproblem Write the standard matrix $A$ for $T$, so that
  $T(x) = Ax$. %
  \subproblem Is $T$ one-to-one?  Either justify why $T$ is
  one-to-one, or write vectors $x$ and $y$ (with $x \neq y$) so that
  $T(x) = T(y)$. %
  \subproblem Is $T$ onto?  Either justify why $T$ is onto or write one
  vector that is not in the range of $T$.
\end{subproblems}

\begin{solution}
  \begin{subproblems}
    \subproblem $A = \mat{T(e_1) T(e_2)}= \mat{1 0; 0 1; 0 1; 1 0}$. %
    \subproblem Yes, $T$ is one-to-one because $A$ has a pivot in
    every column. %
    \subproblem No, $T$ is not onto because $A$ does not have a pivot
    in every row.  From the form of $T$ we see that if $x$ is any
    vector then $T(x)$ has its first and fourth entries equal and also
    has its second and third entries equal. Therefore, any vector
    whose first and fourth entries are different, or whose second and
    third entries are different, will NOT be in the range of $T$.
    Many examples possible, for example $\vec{1 0 0 0}$,
    $\vec{0 1 0 0}$, $\vec{1 2 0 -1}$, $\vec{1 1 0 1}$, etc.
  \end{subproblems}
\end{solution}



\problem % Jankowski quiz 4 Fall 2017
Let $T: {\bf R}^2 \to {\bf R}^2$ be the transformation which reflects about the
line $y = x$, then rotates counterclockwise by $45^\circ$.  Using only geometry,
find $T(e_1)$ and $T(e_2)$.  Use your answer to find the matrix $A$ so that
$T(x)=Ax$ for all $x$ in ${\bf R}^2$.

\begin{solution}
  We examine what happens to $e_1$ and $e_2$.
 \newline
 $e_1$:  The reflection sends $\vec{1 0}$ to $\vec{0 1}$, then rotating cc by $45^\circ$ gives
 $T(e_1) = \mat{\frac{-1}{\sqrt{2}} ; ; \frac{1}{\sqrt{2}} }$.
\newline \newline \newline
$e_2$: The reflection sends $\vec{0 1}$ to $\vec{1 0}$, then rotating cc by $45^\circ$ gives
$T(e_2) = \mat{\frac{1}{\sqrt{2}} ; ; \frac{1}{\sqrt{2}} }$.

Therefore,
\[A = \mat{T(e_1) T(e_2)} = \mat{ \frac{-1}{\sqrt{2}} \frac{1}{\sqrt{2}}; ;
    \frac{1}{\sqrt{2}} \frac{1}{\sqrt{2}}}.\]
\end{solution}

\end{problems} \setcounter{problem}{0}

\newpage

\subsection{3.4, Matrix multiplication}

\begin{problems}
  \problem % Created for ILA
  True or false: If $A$ and $B$ are matrices and $AB$ is an $3 \times 3$ matrix,
  then $A$ and $B$ must be $3 \times 3$ matrices.
  
\begin{solution}
  False: for example, if $A$ is $3 \times 2$ and $B$ is $2 \times 3$, then $AB$
  is $3 \times 3$ but neither $A$ nor $B$ is $3 \times 3$.
\end{solution}

\problem % Created for ILA
True or false: Suppose $A$ and $B$ are matrices so that $AB$ and $BA$ are both
defined.  If $AB$ and $BA$ are each $r \times r$ matrices, then $A$ and $B$ must
each be $n \times n$.

\begin{solution}
  True: $A$ Say $A$ is $k \times \ell$ and $B$ is $m \times n$.  Since $AB$ is
  defined, we know $\ell = m$.  Since $BA$ is defined, we know $n = k$.  This
  means $B$ is $\ell \times k$.  \newline \newline %
  Now: $AB$ is a $k \times \ell$ times an $\ell \times k$, so $AB$ is
  $k \times k$.  Since we were told $AB$ is $r \times r$, this means $r =
  k$. \newline \newline %
  Similarly, $BA$ is an $\ell \times k$ times a $k \times \ell$, so $BA$ is
  $\ell \times \ell$.  Since we were told $BA$ is $r \times r$, this means $r = \ell$.
  \newline \newline %
  Putting it all together, we see $A$ and $B$ are both $r \times r$ matrices.
  
\end{solution}


\problem % Created for ILA
Suppose $A$ is a $24 \times 17$ matrix and the products $AB$ and $BA$ are both
defined.  How many rows does $B$ have?  How many columns does $B$ have?

\begin{solution}
  For $AB$ to be defined, the number of columns of $A$ (namely $17$) must equal
  the number of rows of $B$, so $B$ has $17$ rows. \newline \newline Similarly:
  for $BA$ to be defined, the number of columns of $B$ must equal the number of
  rows of $A$ (namely 24), so $B$ must have $24$ columns. \newline \newline
  Therefore, $B$ is a $17 \times 24$ matrix.
\end{solution}

\problem Suppose $A$ and $B$ are $3 \times 3$ matrices and the columns of $B$
are linearly dependent.  Is it necessarily the columns of $AB$ are linearly
dependent?  Justify your answer.

\begin{solution}
  Yes, the columns of $AB$ are linearly dependent.  Using the fact that $Bv=0$
  for some nonzero $v$ in ${\bf R}^3$, we compute \[ABv=AB(v) = A0=0.\]
  Therefore, the vector $v$ is a non-trivial solution to the system $ABx = 0$,
  so the columns of $AB$ are linearly dependent.
\end{solution}

\problem % Created for ILA
Let $A = \mat{1 -3 4; 1 0 0}$ and $B = \mat{1 \pi}$.  Compute all products that
make sense among $AB$ and $BA$.

\begin{solution}
  $A$ is $2 \times 3$ and $B$ is $1 \times 2$, so the only product that makes
  sense is $BA$.
  \[BA = \mat{1 \pi}\mat{1 -3 4; 1 0 0} = \mat{1+\pi \, -3 4}.\]
\end{solution}

\problem % Created for ILA
Let $T: {\bf R}^2 \to {\bf R}^2$ be the linear transformation that rotates
vectors by $\pi/2$ radians counterclockwise, then scales the $y$-direction by a
factor of $5$, then reflects vectors across the line $y = -x$.  Compute
$T\vec{5 1}$ using matrix multiplication.

\begin{solution}
  The matrix for $T$ is $ABC$ where $C$ is the first operation, $B$ is the
  second operation, and $A$ is the third operation.  Therefore, the matrix for
  $T$ is
  \[\mat{0 -1; -1 0}\mat{1 0; 0 5} \mat{0 -1; 1 0}
    = \mat{0 -1; -1 0} \mat{0 -1; 5 0} = \mat{-5 0; 0 1}.\] %
  We compute
  \[T \vec{5 1} = \mat{-5 0; 0 1}\vec{5 1} = \vec{-25 1}.\] We can check our
  answer geometrically.  $T$ first rotates $\vec{5 1}$ by $\pi/2$ radians
  counterclockwise to get $\vec{-1 5}$, then scales the $y$-coordinate by a
  factor of $5$ to get $\vec{-1 25}$, then reflects across the line $y = -x$ to
  get $\vec{-25 1}$.
\end{solution}


\problem % Created for ILA
$T$ has domain ${\bf R}^{4}$ and codomain ${\bf R}^{7}$. \newline %
$U$ has domain ${\bf R}^{13}$ and codomain ${\bf R}^{25}$. \newline %
$V$ has domain ${\bf R}^{25}$ and codomain ${\bf R}^{17}$. \newline %
$W$ has domain ${\bf R}^{7}$ and codomain ${\bf R}^{13}$. \newline %
Determine which composition of all four transformations makes sense, and find
the number of rows and columns for its standard matrix.

\begin{solution}
  In order for a composition $R \circ S$ to make sense, the domain of $R$ must
  be the codomain of $S$, because we need to be able to take $R(S(v))$ for any
  vector $v$ in the domain of $S$. \newline The composition of all $4$
  transformations that makes sense is $V \circ U \circ W \circ T$.
  \begin{itemize}
  \item If $x$ is a vector in ${\bf R}^{4}$, then $T(x)$ is a vector in
    ${\bf R}^{7}$, which is the domain of $W$. %
  \item Next step: $W(T(x))$ is in ${\bf R}^{13}$, which is the domain of $U$. %
  \item Next step: $U(W(T(x)))$ is in ${\bf R}^{41}$, which is the domain of $V$,
  \item Finally: $V(U(W(T(x))))$ is a vector in ${\bf R}^{17}$.
  \end{itemize}
  Following all of the steps, we see that if $x$ is in ${\bf R}^{4}$ then
  $(V \circ U \circ W \circ T)(x)$ is a vector in ${\bf R}^{17}$, so the
  transformation has domain ${\bf R}^{4}$ and codomain ${\bf R}^{17}$.
  Therefore, its standard matrix is $17 \times 4$, so it has $17$ rows and $4$
  columns.
\end{solution}

\problem % Created for ILA
Let $T$ be the linear transformation $T(x) = ABx$, where
$A$ is $3 \times 3$ and $B$ is $3 \times 2$.
\begin{subproblems}
\subproblem Complete the following: $T$ is onto...
\[\textup{(i) sometimes} \qquad \textup{(ii) always}
  \qquad \textup{(iii) never}\]
\subproblem Complete the following: $T$ is one-to-one...
\[\textup{(i) sometimes} \qquad \textup{(ii) always}
  \qquad \textup{(iii) never}\]
\subproblem Find the formula for $T(x_1, x_2, x_3)$ if
$A = \mat{1 -1 4; 0 1 2; 4 1 0}$ and $B = \mat{1 -1; 2 -2; 0 1}$.
\end{subproblems}

\begin{solution}
  \begin{subproblems}
    \subproblem Never: $AB$ is $3 \times 2$, so it has $3$ rows but a max of $2$
    pivots, therefore at least one row does not have a pivot, thus $T$ cannot be
    one-to-one. %
    \subproblem Sometimes.  $AB$ is $3 \times 2$, so it might have a pivot in
    every column, but this depends on $A$ and $B$.  For example, if
    $B = \mat{0 0; 0 0; 0 0}$ then $AB$ is the $3 \times 2$ zero matrix, so $T$
    is clearly not one-to-one.  \newline If $A = \mat{1 0 0; 0 1 0; 0 0 1}$ and
    $B = \mat{1 0; 0 1; 0 0}$, then $AB = \mat{1 0; 0 1; 0 0}$ which has a pivot
    in every column, so $T$ is one-to-one. %
    \subproblem
    $AB = \mat{1 -1 4; 0 1 2; 4 1 0} \mat{1 -1; 2 -2; 0 1} = %
    \mat{-1 5; 2 2; 6 -6}$.
  \end{subproblems}
\end{solution}


\problem % Jankowski quiz 4 Fall 2017
Let $T: {\bf R}^2 \to {\bf R}^2$ be the transformation which reflects about the
line $y = x$, then rotates counterclockwise by $45^\circ$.
\begin{subproblems}
  \subproblem Write the matrix $J$ for reflection across the line $y=x$ in
  ${\bf R}^2$.  %
  \subproblem Write the matrix $K$ for rotation counterclockwise
  by $45^\circ$ in ${\bf R}^2$. %
  \subproblem Using matrix multiplication,
  compute the matrix $A$ so that $T(x)=Ax$ for all $x$ in ${\bf R}^2$.
\end{subproblems}

\begin{solution}
\begin{subproblems}
  \subproblem $J=\mat{0 1; ; 1 0}$.
  \subproblem $K=\mat{ \frac{1}{\sqrt{2}} \frac{-1}{\sqrt{2}}; ;
    \frac{1}{\sqrt{2}} \frac{1}{\sqrt{2}}}$.
  \subproblem $A=KJ$:
  \[ A = \mat{ \frac{1}{\sqrt{2}} \frac{-1}{\sqrt{2}}; ; \frac{1}{\sqrt{2}} \frac{1}{\sqrt{2}}}
  \mat{0 1; ; 1 0}
  = \mat{ \frac{-1}{\sqrt{2}} \frac{1}{\sqrt{2}}; ; \frac{1}{\sqrt{2}} \frac{1}{\sqrt{2}}}.\]
\end{subproblems}
\end{solution}


\problem % Jankowski Spring 2018 exam 2
Let $T: {\bf R}^2 \to {\bf R}^2$ be the transformation of reflection
about the line $y= -x$, and let $U: {\bf R}^2 \to {\bf R}^3$ be the
transformation $U \vec{x y} = \vec{x-y x 3y}$.
\begin{subproblems}
  \subproblem Write the standard matrix $A$ for $T$.  Is $T$ onto? %
  \subproblem Write the standard matrix $B$ for $U$.  Is $U$
  one-to-one?  Briefly justify your answer. %
  \subproblem Circle the composition that makes sense: \quad
  $T \circ U$ \qquad $U \circ T$ %
  \subproblem Compute the standard matrix for the composition you
  circled in part (c).
\end{subproblems}

\begin{solution}
  \begin{subproblems}
    \subproblem $A=\mat{0 -1; -1 0}$. Yes, $T$ is onto.  One way to
    see this is to note that $A$ has a pivot in every row.
    \subproblem $B=\mat{U(e_1) U(e_2)} = \mat{1 -1; 1 0; 0 3}$.  Yes,
    $U$ is one-to-one because $B$ has a pivot in every column.
    \subproblem $U \circ T$ makes sense since it sends
    $R^2 \to R^2 \to R^3$.  \subproblem
    $BA = \mat{1 -1; 1 0; 0 3} \mat{0 -1; -1 0} =
    \mat{1 -1; 0 -1; -3 0}$.
  \end{subproblems}
\end{solution}

\problem Consider the following transformations. \newline %
$T: {\bf R}^2 \to {\bf R}^2$ be rotation by $90^\circ$ clockwise.
\newline %
$U: {\bf R}^3 \to {\bf R}^2$ given by $U(a, b, c) = (a-b, c-5a)$.
\newline %
$V: {\bf R}^{4} \to {\bf R}^3$ given by
$V(x, y, z, w) = (w, x-y+z, z-x)$.
\begin{subproblems}
  \subproblem Which composition of all three transformation makes
  sense? %
  \subproblem Using matrix multiplication, compute the standard
  matrix for the transformation you wrote in part (a).
\end{subproblems}

\begin{solution}
  \begin{subproblems}
    \subproblem $T \circ U \circ V$ makes sense: $V$ sends a vector
    $x$ from ${\bf R}^4$ to ${\bf R}^3$, then $U$ sends $V(x)$ from
    ${\bf R}^3$ to ${\bf R}^2$, then $T$ sends $U(V(x))$ from
    ${\bf R}^2$ to ${\bf R}^2$. %
    \subproblem $T(x) = Ax$, $U(x) = Bx$, and $V(x) = Cx$ for the
    matrices below, where we multiply $ABC$ to get the matrix for
    $T \circ U \circ V$.  Matrix multiplication is NOT commutative,
    but it is associative, so to compute $ABC$ we can either compute
    $A(BC)$ or $(AB)C$.
    \begin{align*}
      ABC 
      &
        = \mat{0 1; -1 0}
        \mat{1 -1 0; -5 0 1}
        \mat{0 0 0 1; 1 -1 1 0; -1 0 1 0} \\
      &
        = \mat{0 1; -1 0}
        \mat{-1 1 -1 1; -1 0 1 -5} \\
      &
        = \mat{-1 0 1 5; 1 -1 1 -1}.
    \end{align*}
  \end{subproblems}
\end{solution}


\end{problems} \setcounter{problem}{0}

\newpage

\subsection{3.5 and 3.6, Matrix inverses and the Invertible Matrix Theorem}

\begin{problems}
%a
%  
\problem % Made for ILA
True or false: Suppose $A$ is an $n \times n$ matrix and $AB = I_n$. 
Then $BA = I_n$.  \newline %
\begin{solution}
 True.  Even though matrix multiplication is generally
    not commutative, it is a general fact that if two $n \times n$
    matrices $A$ and $B$ satisfy $AB=I$ then $BA=I$, and similarly if
    $BA=I$ then $AB=I$. %   
\end{solution}

\problem % Made for ILA
True or false: There is an invertible matrix $A$ satisfying 
$A \vec{1 2 3} = \vec{1 2}$. 

\begin{solution}
False.  Only square matrices cann be invertible.  If
    $A \vec{1 2 3} = \vec{1 2}$ then $A$ is $2 \times 3$, so $A$
    cannot be invertible. %   
\end{solution}

% Jankowski Spring 2018 exam 2 for 3.6
\problem True or false: If $A$ is an $n \times n$ matrix and $Ax=0$ has 
only the trivial solution, then the equation $Ax=b$ is consistent for 
every $b$ in ${\bf R}^n$. 

\begin{solution}
  True.  $A$ is invertible by the Invertible Matrix Theorem.   
\end{solution}

% Jankowski Spring 2018 exam 2 for 3.6
\problem True or false: There are linear transformations $T: {\bf R}^3 
\to {\bf R}^4$ and $U: {\bf R}^4 \to {\bf R}^3$ so that $T \circ U$ is 
invertible.

\begin{solution}
False.  Since $T$ and $U$ are linear, they are matrix transformations $T(v) = Av$ and $U(x) = Bx$ for some matrices $A$ and $B$.  Since $U$ has domain ${\bf R}^4$ and codomain ${\bf R}^3$, we know that $B$ is $3 \times 4$, so $B$ has $4$ columns but at most $3$ pivots.  Therefore, $B$ has at least one column without a pivot, so $Bx = 0$ has infinitely many solutions.  For any solution to $Bx=0$, we get
\[(T \circ U)(x) = ABx = A(Bx)=A(0) = 0,\]
so $(T \circ U)(x)=0$ has infinitely many solutions, thus $T \circ U$ cannot be one-to-one.  By the Invertible Matrix Theorem, $T \circ U$ cannot be invertible.
\end{solution}

% Jankowski Spring 2018 exam 2 for 3.6
\problem True or false: There are linear transformations $T: {\bf R}^4 
\to {\bf R}^3$ and $U: {\bf R}^3 \to {\bf R}^4$ so that $T \circ U$ is 
invertible.

\begin{solution}
  True.  Take $U(x_1, x_2, x_3)= (x_1, x_2, x_3, 0)$ and
  $T(x_1, x_2, x_3, x_4)= (x_1, x_2, x_3)$.  Then
  $(T \circ U)(x_1, x_2, x_3)= (x_1, x_2, x_3)$ so $T \circ U$ is invertible (in
  fact, it is the identity transformation!).
\end{solution}


    
\problem % Made for ILA
True or false: Let $A = \mat{2 0 0; 0 3 0; 0 0 1}$. Then $A$ is 
invertible and $A^{-1} = \mat{\frac{1}{2} 0 0; 0 \frac{1}{3} 0; 0 0 1}$. 

\begin{solution}
True:
    $\mat{2 0 0; 0 3 0; 0 0 1} \mat{\frac{1}{2} 0 0; 0 \frac{1}{3} 0;
      0 0 1} = \mat{1 0 0; 0 1 0; 0 0 1}=I.$ \newline This shows that
    the inverse of $A$ is
    $\mat{\frac{1}{2} 0 0; 0 \frac{1}{3} 0; 0 0 1}$. %    
\end{solution}

\problem % Made for ILA
True or false: If $A$ is an $n \times n$ matrix and $Ax = 0$ has 
infinitely many solutions, then $A$ cannot be invertible.

\begin{solution}
True.  If $A$ is invertible then for each $b$ in
    ${\bf R}^n$, the equation $Ax=b$ has a unique solution
    $x = A^{-1}b$.  In particular the homogeneous equation $Ax=0$ must
    have unique solution $x= A^{-1}0 = 0$.   
\end{solution}

\problem % Made for ILA
Suppose $A$ is the matrix whose inverse is given by
$A^{-1} = \mat{-4 1 1; 0 2 3; -1 2 2}$.  \newline %
Find all solutions to $Ax = \vec{1 -3 9}$. \newline %
\begin{solution}
Since $A$ is invertible, any matrix equation $Ax = b$ (for any $b$ in ${\bf R}^3$) has the unique solution $x = A^{-1} b$.  
\[Ax = \vec{1 -3 9}, \qquad x = A^{-1}\vec{1 -3 9}
= \mat{-4 1 1; 0 2 3; -1 2 2} \vec{1 -3 9} = 
\vec{2 21 11}.\]
% In reality, $A = \dfrac{1}{7} \mat{-2 0 1; -3 -7 12; 2 7 -8}$
\end{solution}

\problem % Made for ILA
Let $A = \mat{7 8; 1 3}$. Is $A$ invertible?  If so, find $A^{-1}$.
\newline %
\begin{solution}
If $A = \mat{a b; c d}$ with $ad - bc \neq 0$, then $A$ is invertible and
$A^{-1} = \dfrac{1}{ad-bc}\mat{d -b; -c a}$. \newline
Here $A=\mat{7 8; 1 3}$ where $7(3)-8(1) = 13 \neq 0$, so $A$ is 
invertible and 
\[A^{-1} = \dfrac{1}{13} \mat{3 -8; -1 7}.\]
\end{solution}



\problem % Made for ILA
Let $A = \mat{4 1 1; 2 6 0; 3 -1 1}$.  Find $A^{-1}$.

\begin{solution}
  We put $\amat{A, I}$ into RREF.
  \begin{align*}
    \mat{4 1 1 | 1 0 0; 2 6 0 | 0 1 0; 3 -1 1 | 0 0 1}
    & \xrightarrow[\textup{then } R_1 = R_1/2]{R_1 \leftrightarrow R_2}
    \mat{1 3 0 | 0 \frac{1}{2} 0; 4 1 1 | 1 0 0; 3 -1 1 | 0 0 1}
    \\ & \xrightarrow[R_3 = R_3-3R_1]{R_2 = R_2 - 4R_1}
         \mat{1 3 0 | 0 \frac{1}{2} 0; 0 -11 1 | 1 -2 0; 0 -10 1
         | 0 -\frac{3}{2} 1}
    \\ & \xrightarrow[]{R_2 = R_2/-11}
    \mat{1 3 0 | 0 \frac{1}{2} 0; 0 1 -\frac{1}{11} |
      -\frac{1}{11} \frac{2}{11} 0; 0 -10 1 | 0 -\frac{3}{2} 1}
    \\ & \xrightarrow[]{R_3 = R_3 + 10R_2}
    \mat{1 3 0 | 0 \frac{1}{2} 0; 0 1 -\frac{1}{11} |
         -\frac{1}{11} \frac{2}{11} 0; 0 0 \frac{1}{11}
         | -\frac{10}{11} \frac{7}{22} 1}
    \\ & \xrightarrow[\textup{then } R_3 = 11R_3]{R_2 = R_2 + R_3}
    \mat{1 3 0 | 0 \frac{1}{2} 0; 0 1 0 |
      -1 \frac{1}{2} 1; 0 0 1 | -10 \frac{7}{2} 11}
    \\ & \xrightarrow[]{R_1 = R_1 - 3R_2}
         \mat{1 0 0 | 3 -1 -3; 0 1 0 | -1 \frac{1}{2} 1; 0 0 1
         | -10 \frac{7}{2} 11}.
   \end{align*}
   Therefore,
   \[ A^{-1} = \mat{1 0 0 | 3 -1 -3; 0 1 0 | -1 \frac{1}{2} 1; 0 0 1 |
       -10 \frac{7}{2} 11}.\]
\end{solution}


\problem % Jankowski Spring 2018 exam 2
Find the matrix $A$ satisfying $A^{-1} e_1 = \vec{0 1 2}$, \
$A^{-1} e_2 = \vec{1 0 0}$, and \ $A^{-1}e_3 = \vec{0 1 3}$.

\begin{solution}
  \begin{align*} \mat{A^{-1} | I} & = \mat{0 1 0 | 1 0 0; 1 0 1 | 0 1
      0; 2 0 3 | 0 0 1} \xrightarrow[]{R_1 \leftrightarrow R_2} \mat{1
      0 1 | 0 1 0; 0 1 0 | 1 0 0; 2 0 3 | 0 0 1} \\ \\ &
    \xrightarrow[]{R_3=R_3-2R_1} \mat{1 0 1 | 0 1 0; 0 1 0 | 1 0 0; 0
      0 1 | 0 -2 1} \xrightarrow[]{R_1=R_1-R_3} \mat{1 0 0 | 0 3 -1; 0
      1 0 | 1 0 0; 0 0 1 | 0 -2 1}.
  \end{align*}  
  Therefore, $A=(A^{-1})^{-1} = \mat{0 3 -1; 1 0 0; 0 -2 1}.$
\end{solution}


\end{problems}\setcounter{problem}{0}


\newpage